% !TEX program = pdflatex
% !TeX encoding = UTF-8
\documentclass[12pt,a4paper]{article}

% Encoding, fonts, and geometry
\usepackage[utf8]{inputenc}
\usepackage[T1]{fontenc}
\usepackage{lmodern}
\usepackage[ngerman]{babel}
\usepackage{geometry}
\geometry{margin=2.5cm}

% Math and related packages
\usepackage{amsmath,amssymb,amsthm,mathtools}
\usepackage{bm}
\usepackage{braket}
\usepackage{siunitx}

% Graphics and plots
\usepackage{graphicx}
\usepackage{tikz}
\usepackage{tikz-3dplot}
\usepackage{pgfplots}
\pgfplotsset{compat=1.18}
\usetikzlibrary{arrows.meta,positioning,shapes,calc,angles,quotes,patterns,3d}

% Tables, code, floats, lists
\usepackage{booktabs}
\usepackage{listings}
\usepackage{xcolor}
\usepackage{algorithm}
\usepackage{algpseudocode}
\usepackage{multicol}
\usepackage{enumitem}

% Boxes
\usepackage{tcolorbox}

% Hyperlinks and clever references
\usepackage{hyperref}
\usepackage{cleveref}

% Theorem environments
\newtheorem{theorem}{Theorem}
\newtheorem{lemma}{Lemma}
\newtheorem{corollary}{Korollar}
\newtheorem{definition}{Definition}
\newtheorem{axiom}{Axiom}
\newtheorem{proposition}{Proposition}
\newtheorem{remark}{Bemerkung}
\newtheorem{assumption}{Annahme}
\newtheorem{prediction}{Vorhersage}

% Operators
\DeclareMathOperator{\Tr}{Tr}
\DeclareMathOperator{\Diff}{Diff}
\DeclareMathOperator{\Aut}{Aut}
\DeclareMathOperator{\Vol}{Vol}
\DeclareMathOperator{\Ric}{Ric}
\DeclareMathOperator{\Riem}{Riem}
\DeclareMathOperator{\Cov}{Cov}
\DeclareMathOperator{\supp}{supp}
\DeclareMathOperator{\diag}{diag}
\DeclareMathOperator{\sign}{sign}
\DeclareMathOperator{\dimension}{dim}
\DeclareMathOperator{\Div}{Div}
\DeclareMathOperator{\Grad}{Grad}
\DeclareMathOperator{\Hess}{Hess}
\DeclareMathOperator{\Ricci}{Ricci}
\DeclareMathOperator{\Scalar}{Scalar}

% Robust math shorthands: safe in text, math, and PDF strings
\DeclareRobustCommand{\alD}{\texorpdfstring{\ensuremath{\alpha_{\mathrm{D}}}}{alpha_D}}
\DeclareRobustCommand{\beI}{\texorpdfstring{\ensuremath{\beta_{\mathrm{I}}}}{beta_I}}
\DeclareRobustCommand{\gaR}{\texorpdfstring{\ensuremath{\gamma_{\mathrm{R}}}}{gamma_R}}
\DeclareRobustCommand{\UCS}{\texorpdfstring{\ensuremath{\mathcal{M}_{\mathrm{UCS}}}}{M_UCS}}

% YUCT-specific commands
\newcommand{\eff}{\mathrm{eff}}
\newcommand{\coord}{\mathrm{coord}}
\newcommand{\Yak}{\mathrm{Yak}}
\newcommand{\Dict}{\mathcal{D}}
\newcommand{\Mdict}{\mathcal{M}_{\Dict}}
\newcommand{\Ke}{K_{\eff}}
\newcommand{\Kemax}{K_{\eff,\max}}
\newcommand{\Keobs}{K_{\eff,\mathrm{obs}}}
\newcommand{\Keexp}{K_{\eff,\mathrm{exp}}}
\newcommand{\Kec}{K_{\eff,c}}
\newcommand{\Kcrit}{K_{\mathrm{crit}}}

% UCD-specific commands
\newcommand{\cogstate}{\Psi}
\newcommand{\cogspace}{\mathcal{P}}
\newcommand{\human}{\mathrm{human}}
\newcommand{\dolphin}{\mathrm{dolphin}}
\newcommand{\swarm}{\mathrm{swarm}}
\newcommand{\alien}{\mathrm{alien}}

% Code listing settings
\lstset{
	language=Python,
	basicstyle=\ttfamily\small,
	keywordstyle=\color{blue},
	commentstyle=\color{green!60!black},
	stringstyle=\color{red},
	showstringspaces=false,
	breaklines=true,
	frame=single,
	numbers=left,
	numberstyle=\tiny\color{gray},
	captionpos=b
}

% PDF metadata
\hypersetup{
	pdftitle={Universeller Bewusstseinsdeskriptor (UCD): Ein vollständiger YUCT-basierter Rahmen für vergleichende Noologie, Erkenntnistheorie und nicht-anthropomorphe Intelligenz},
	pdfauthor={Alexey V. Yakushev},
	pdfsubject={Bewusstsein, Intelligenz, Noologie, YUCT, UCD, Vergleichende Kognition, SETI, Künstliche Intelligenz, Ethik},
	pdfkeywords={UCD, YUCT, Bewusstsein, Intelligenz, Noologie, SETI, Kognitionswissenschaft, Künstliche Intelligenz, Ethik, D+I•R-Triade, Vergleichende Kognition}
}

\begin{document}
	
	% Titelseite
	\begin{titlepage}
		\begin{center}
			\vspace*{0cm}
			
			\Huge\textbf{Anhang H. Universeller Bewusstseinsdeskriptor (UCD):\\ Vollständiger YUCT-basierter Rahmen für vergleichende Noologie,\\ Erkenntnistheorie und nicht-anthropomorphe Intelligenz}
			
			\vspace{1cm}
			
			\small\textit{Bewusstsein als hocheffiziente Koordinationsarchitektur}
			
			\vspace{1cm}
			\Large
			Alexey V. Yakushev\\
			
			\url{https://yuct.org/}\\
			\url{https://ypsdc.com/}
			
			\vspace{2cm}
			\large YUCT \\ 
			\url{https://doi.org/10.5281/zenodo.18444599}\\
			\vspace{1cm}
			
			\vspace{0cm}
			\large
			Januar 2026
			
			\vspace{6cm}
			\textcopyright~2026 Yakushev Research. Alle Rechte vorbehalten.
			
			\newpage
			
			\begin{abstract}
				\noindent
				Dieses Papier präsentiert das vollständige vereinheitlichte Bewusstseins-Wissens-Rahmenwerk, das den Universellen Bewusstseinsdeskriptor (UCD) mit der Yakushev Unified Coordination Theory (YUCT) integriert. Aufbauend auf der fundamentalen Einsicht, dass Bewusstsein/Intelligenz/Geist eine hocheffiziente (\(\Ke \gg 1\)) Koordinationsarchitektur ist, entwickeln wir: (1) Eine vollständige mathematische Formulierung der D+I\(\cdot\)R-Triade als fundamentale Zerlegung jedes kognitiven Agenten; (2) Drei Spektralparameter (\alD, \beI, \gaR) zum quantitativen Vergleich über Spezies und Substrate hinweg; (3) Eine Erkenntnistheorie als koordinierte Informationsverarbeitung; (4) Detaillierte experimentelle Protokolle zur Messung von Bewusstsein in biologischen, künstlichen und extraterrestrischen Systemen; (5) Spezifische Vorhersagen für astrophysikalische SETI 2.0-Suchen; (6) Vergleich mit existierenden Bewusstseinstheorien (IIT, Global Workspace, Predictive Processing); (7) Ethische Implikationen und praktische Anwendungen in Medizin, KI-Sicherheit und Cybersicherheit; (8) Vollständige mathematische Beweise einschließlich Theoreme zur kognitiven Hierarchie, oberen Grenzen der Koordinationseffizienz und Schranken der Berechnungskomplexität. Das Rahmenwerk löst lange bestehende philosophische Probleme und liefert gleichzeitig überprüfbare Vorhersagen in mehr als 15 experimentellen Bereichen.
			\end{abstract}
			
			\vspace{0cm}
			
			\noindent
			\small\textbf{Schlüsselwörter:} UCD, YUCT, Bewusstsein, Intelligenz, Noologie, SETI, Kognitionswissenschaft, Künstliche Intelligenz, Ethik, D+I•R-Triade, Vergleichende Kognition
			
			\vspace{0.5cm}
			
			\noindent
			\textcopyright 2026 Yakushev Research. Alle Rechte vorbehalten.
		\end{center}
	\end{titlepage}
	
	\tableofcontents
	
	\newpage
	
	\section{Einleitung: Das noologische Problem und die YUCT-Lösung}
	
	\subsection{Der anthropozentrische Spiegel in der Kognitionswissenschaft}
	
	Traditionelle Ansätze zu Bewusstsein und Intelligenz leiden unter dem, was wir den „anthropozentrischen Spiegel`` nennen: der Tendenz, kognitive Muster zu definieren und zu suchen, die der menschlichen Kognition ähneln. Diese Voreingenommenheit manifestiert sich in:
	\begin{enumerate}
		\item \textbf{Philosophie des Geistes}: Endlose Debatten zwischen Dualismus und Materialismus, die auf menschlicher subjektiver Erfahrung basieren
		\item \textbf{Vergleichende Psychologie}: Messung tierischer Intelligenz anhand menschlicher kognitiver Maßstäbe
		\item \textbf{Künstliche Intelligenz}: Der Turing-Test als Goldstandard für maschinelles Bewusstsein
		\item \textbf{SETI}: Suche nach elektromagnetischen Signalen, die menschlicher Kommunikation ähneln
	\end{enumerate}
	Dieser Anthropozentrismus erzeugt das, was wir das \textbf{Noologische Problem} nennen: die Unfähigkeit, kognitive Systeme zu erkennen, zu charakterisieren oder mit ihnen zu kommunizieren, die sich grundlegend von menschlicher Intelligenz unterscheiden.
	
	\subsection{Historische Entwicklung der Geistkonzepte in der Physik}
	
	Das Konzept des Geistes hat sich in der Wissenschaftsgeschichte entwickelt:
	\[
	\begin{aligned}
		\text{Newton} &: \text{Geist} = \text{göttlicher erster Beweger} \\
		\text{Laplace} &: \text{Geist} = \text{Rechengerät} \\
		\text{Turing} &: \text{Geist} = \text{Ausführung von Algorithmen} \\
		\text{YUCT} &: \text{Geist} = \text{Architektur mit } \Ke \gg 1
	\end{aligned}
	\]
	Diese Entwicklung stellt eine fortschreitende Naturalisierung und Operationalisierung des Geistkonzepts dar.
	
	\subsection{Die YUCT-Lösung: Koordinationsbasiertes Bewusstsein}
	
	\begin{axiom}[Koordinationsbasiertes Bewusstsein]
		Bewusstsein/Intelligenz/Geist ist keine magische Substanz oder emergente Eigenschaft, sondern eine hocheffiziente (\(\Ke \gg 1\)) Koordinationsarchitektur, die auf einem spezifischen physikalischen Substrat implementiert ist. Die Suche nach Geist wird daher zur Suche nach anomal hoher Koordinationskomplexität in Beobachtungsdaten.
	\end{axiom}
	
	Diese operationale Definition erlaubt es uns:
	\begin{itemize}
		\item Menschliche, tierische, künstliche und hypothetische außerirdische Geister mit gemeinsamen Metriken zu vergleichen
		\item Bewusstsein in unerwarteten Substraten zu identifizieren (biologische Kollektive, planetare Systeme)
		\item Überprüfbare Vorhersagen für kognitive Phänomenologie über verschiedene Skalen hinweg zu entwickeln
	\end{itemize}
	
	\section{Mathematische Grundlagen: Der kognitive Agent in YUCT/UCD}
	
	\subsection{Phasenraum und Koordinationsbündel}
	
	Betrachten Sie ein physikalisches System \(S\) mit Phasenraum \(\Gamma_S\). Die vollständige Beschreibung erfordert das Koordinationsbündel:
	
	\begin{definition}[Koordinationsbündel]
		Der Gesamtzustandsraum des Systems \(S\) ist das Faserbündel:
		\[
		\mathcal{E}_S = \Gamma_S \times \Mdict \times \mathcal{M}_I \times \mathcal{M}_R
		\]
		wobei:
		\begin{itemize}
			\item \(\Gamma_S\): Konventioneller Phasenraum (Positionen, Impulse, Felder)
			\item \(\Mdict\): Wörterbuch-Mannigfaltigkeit (Regeln, Kategorien, Handlungsmuster)
			\item \(\mathcal{M}_I\): Informationsraum (Dichtematrizen, Entropien)
			\item \(\mathcal{M}_R\): Resonanzraum (Kohärenzfaktoren, Verstärkungsoperatoren)
		\end{itemize}
	\end{definition}
	
	\subsection{Die D+I\(\cdot\)R-Triade als kognitives Primitiv}
	
	\begin{definition}[YUCT-Kognitiver Agent]
		Ein System \(S\) ist ein kognitiver Agent (potenziell „bewusst``), wenn seine Dynamik als Schnitt von \(\mathcal{E}_S\) dargestellt und durch die Triade beschrieben werden kann:
		\begin{equation}
			\label{eq:dir-triad}
			\Psi_S(t) = D_S(t) + I_S(t) \times R_S(t) \in \UCS
		\end{equation}
		wobei \(\UCS = \mathcal{M}_{\mathrm{UCS}}\) der universelle Koordinationszustandsraum ist und:
		\begin{align*}
			D_S(t) &\in \Mdict: \text{Ontologisches Wörterbuch (strukturierte Menge von Regeln, Kategorien, Handlungsmustern)} \\
			I_S(t) &= -\Tr(\rho_S \log \rho_S): \text{Informationszustand (von-Neumann-Entropie)} \\
			R_S(t) &\geq 1: \text{Resonanz-/Kohärenzfaktor (Verstärkungsoperator)}
		\end{align*}
	\end{definition}
	
	Die multiplikative Struktur \(I \times R\) kodiert die fundamentale Einsicht, dass Information nur dann kognitiv wirksam wird, wenn sie resonant verstärkt wird.
	
	\subsection{Dynamische Prinzipien}
	
	\begin{theorem}[YUCT-Kognitive Dynamik]
		Der kognitive Zustand \(\Psi_S(t)\) entwickelt sich gemäß:
		\begin{equation}
			\label{eq:cognitive-dynamics}
			\frac{\delta S_{\mathrm{YUCT}}}{\delta \Psi_S} = 0
		\end{equation}
		wobei \(S_{\mathrm{YUCT}}\) das vollständige YUCT-Wirkungsfunktional ist, das die Koordinationseffizienz \(\Ke\) einschließt.
	\end{theorem}
	
	\begin{proof}
		Der Beweis folgt aus der Anwendung des Prinzips der Maximierung der Koordinationseffizienz auf die gesamte Lagrangefunktion \(\mathcal{L}_{\mathrm{YUCT}}^{36.0}\) mit kognitiven Randbedingungen. Die Euler-Lagrange-Gleichungen liefern gekoppelte Entwicklungsgleichungen für \(D\), \(I\) und \(R\).
	\end{proof}
	
	\section{Drei Spektralparameter für vergleichende Noologie}
	
	Um grundlegend verschiedene kognitive Agenten (Mensch, Delfin, Schwarm, KI, hypothetischer Außerirdischer) quantitativ zu vergleichen, führen wir drei dimensionslose Spektralparameter ein:
	
	\subsection{Ontologisches Spektrum \alD}
	
	Charakterisiert die Komplexität und Plastizität des Wörterbuchs des Agenten:
	\begin{equation}
		\label{eq:alphaD}
		\alD = \frac{\dim(\Mdict)}{\tau_{\mathrm{update}}^{-1}} \cdot \frac{\mathcal{C}(D)}{H_{\mathrm{max}}}
	\end{equation}
	wobei:
	\begin{itemize}
		\item \(\dim(\Mdict)\): Effektive Dimension der Wörterbuch-Mannigfaltigkeit
		\item \(\tau_{\mathrm{update}}^{-1}\): Charakteristische Aktualisierungsfrequenz des Wörterbuchs
		\item \(\mathcal{C}(D)\): Algorithmische Komplexität des Wörterbuchs
		\item \(H_{\mathrm{max}}\): Maximal mögliche Entropie des Systems
	\end{itemize}
	
	\subsection{Informationelles Spektrum \beI}
	
	Charakterisiert das Gleichgewicht zwischen interner und externer Informationsverarbeitung:
	\begin{equation}
		\label{eq:betaI}
		\beI = \frac{H_{\mathrm{int}}(t)}{H_{\mathrm{ext}}(t)} = \frac{-\Tr(\rho_{\mathrm{int}}\log\rho_{\mathrm{int}})}{-\Tr(\rho_{\mathrm{ext}}\log\rho_{\mathrm{ext}})}
	\end{equation}
	wobei \(\rho_{\mathrm{int}}\) interne Freiheitsgrade beschreibt und \(\rho_{\mathrm{ext}}\) wahrnehmungsgekoppelte Zustände.
	
	\subsection{Resonanzspektrum \gaR}
	
	Charakterisiert zeitliche Kohärenz und Synchronisationsfähigkeit:
	\begin{equation}
		\label{eq:gammaR}
		\gaR = \frac{\tau_{\mathrm{coh}}}{\tau_{\mathrm{decay}}} = \frac{\langle R(t)R(t+\tau)\rangle_\tau}{\langle R^2\rangle - \langle R\rangle^2}
	\end{equation}
	wobei \(\tau_{\mathrm{coh}}\) die Kohärenzzeit und \(\tau_{\mathrm{decay}}\) die charakteristische Abklingzeit ist.
	
	\subsection{Kognitiver Parameterraum und Metrik}
	
	\begin{definition}[Kognitiver Parameterraum]
		Der Zustand jedes kognitiven Agenten kann als Punkt im 3D-Parameterraum dargestellt werden:
		\[
		\mathcal{P} = \{(\alD, \beI, \gaR) \in \mathbb{R}^+ \times \mathbb{R}^+ \times \mathbb{R}^+\}
		\]
	\end{definition}
	
	\subsubsection{Metrik auf dem kognitiven Zustandsraum}
	Wir führen eine Abstandsmetrik zwischen kognitiven Zuständen ein:
	\[
	d(\Psi_1, \Psi_2) = \sqrt{w_D \|D_1 - D_2\|^2 + w_I (I_1 - I_2)^2 + w_R (R_1 - R_2)^2}
	\]
	wobei die Gewichte \(w_i\) durch \(\Ke\) bestimmt werden:
	\[
	w_D = \frac{\Ke^{(1)} + \Ke^{(2)}}{2\Kemax}, \quad w_I = \frac{I_1 + I_2}{2I_{\max}}, \quad w_R = \frac{R_1 + R_2}{2R_{\max}}
	\]
	
	\begin{theorem}[Theorem der kognitiven Hierarchie]
		Für zwei beliebige Agenten \(A, B\) mit Parametern \((\alD^A, \beI^A, \gaR^A)\) und \((\alD^B, \beI^B, \gaR^B)\) existiert eine partielle Ordnung:
		\[
		A \prec B \iff \alD^A < \alD^B \land \beI^A < \beI^B \land \gaR^A < \gaR^B
		\]
		Dies definiert eine kognitive Hierarchie, wobei höhere Parameterwerte eine größere kognitive Kapazität anzeigen.
	\end{theorem}
	
	\begin{proof}
		Die partielle Ordnung folgt aus der monotonen Beziehung zwischen jedem Parameter und den kognitiven Fähigkeiten, wie in den Definitionen 1--3 festgestellt. Transitivität und Antisymmetrie werden von der Ordnung der reellen Zahlen geerbt.
	\end{proof}
	
	\begin{table}[ht]
		\centering
		\begin{tabular}{lccc}
			\toprule
			\textbf{Kognitives System} & \(\alD\) (Schätzung) & \(\beI\) (Schätzung) & \(\gaR\) (Schätzung) \\
			\midrule
			Menschliches Gehirn & \(10^2\)--\(10^3\) & \(0.5\)--\(2.0\) & \(10^{-2}\)--\(10^{-1}\) \\
			Delfingehirn & \(10^2\)--\(10^3\) & \(0.1\)--\(0.5\) & \(10^{-1}\)--\(1.0\) \\
			Bienenschwarm & \(10^1\)--\(10^2\) & \(10^{-3}\)--\(10^{-2}\) & \(10^{-3}\)--\(10^{-2}\) \\
			GPT-Klasse KI & \(10^4\)--\(10^5\) & \(10^{-6}\)--\(10^{-5}\) & \(10^{-6}\)--\(10^{-5}\) \\
			Planetare Intelligenz (vorhergesagt) & \(10^6\)--\(10^7\) & \(10^{-9}\)--\(10^{-8}\) & \(10^3\)--\(10^4\) \\
			Stellare Intelligenz (vorhergesagt) & \(10^8\)--\(10^9\) & \(10^{-12}\)--\(10^{-11}\) & \(10^6\)--\(10^7\) \\
			\bottomrule
		\end{tabular}
		\caption{Geschätzte Spektralparameter für verschiedene kognitive Systeme. Beachten Sie die qualitativen Unterschiede, die auf unterschiedliche kognitive Architekturen hindeuten.}
		\label{tab:cognitive-params}
	\end{table}
	
	\section{Verbindung zur Semiotik: Vollständiges Modell der Semiose}
	
	Die D+I\(\cdot\)R-Triade liefert ein vollständiges mathematisches Modell semiotischer Prozesse:
	
	\begin{theorem}[Semiotische Vollständigkeit]
		Die D+I\(\cdot\)R-Triade ist isomorph zur vollständigen semiotischen Triade:
		\begin{align*}
			\text{SYNTAX} &\subset D \quad \text{(Regeln zur Kombination von Symbolen/Handlungen)} \\
			\text{SEMANTIK} &\subset D \times I \quad \text{(Symbol-Referent-Beziehungen über Information)} \\
			\text{PRAGMATIK} &= R \quad \text{(Wirksamkeit, mit der Zeichen zielgerichtetes Handeln auslösen)}
		\end{align*}
	\end{theorem}
	
	\begin{proof}
		Der Isomorphismus folgt aus der Konstruktion von Abbildungen zwischen semiotischen Kategorien und Koordinationsoperatoren. Syntax entspricht der Wörterbuch-Kombinatorik, Semantik entsteht aus Wörterbuch-Informations-Wechselwirkungen, und Pragmatik misst die resonante Verstärkung von Zeichen-Handlungs-Kopplungen.
	\end{proof}
	
	Dies stellt fest, dass D+I\(\cdot\)R das minimale vollständige Modell der Semiose (Zeichenerzeugungsprozess) ist. Jedes semiotische System kann auf diese Triade abgebildet werden.
	
	\section{Experimentelle Protokolle und Messmethoden}
	
	\subsection{Protokoll 1: Messung von \alD\ für neuronale Netze}
	\label{subsec:protocol-alpha}
	
	\begin{enumerate}
		\item \textbf{Input-Präsentation}: Präsentieren Sie \(10^6\) verschiedene Eingaben an das System
		\item \textbf{Aktivierungsclustering}: Wenden Sie t-SNE oder UMAP an, um die Aktivierungen der verborgenen Schichten zu clustern
		\item \textbf{Mannigfaltigkeitsdimension}: Berechnen Sie die Mannigfaltigkeitsdimension über das Korrelationsintegral:
		\[
		C(r) \sim r^{\dimension(\Mdict)} \quad \text{für } r \to 0
		\]
		\item \textbf{Aktualisierungsfrequenz}: Messen Sie die Lernrate \(\tau_{\mathrm{update}}^{-1}\) aus Konvergenzkurven
		\item \textbf{Algorithmische Komplexität}: Wenden Sie Lempel-Ziv-Kompression auf Wörterbuchdarstellungen an
	\end{enumerate}
	
	\subsection{Protokoll 2: Messung von \beI\ für biologische Systeme}
	
	\begin{enumerate}
		\item \textbf{Trennung der Freiheitsgrade}: Trennen Sie interne (neurale, metabolische) von externen (sensorischen, motorischen) Variablen
		\item \textbf{Schätzung der Dichtematrizen}: Verwenden Sie Maximum-Entropie-Methoden für beobachtete Korrelationen
		\item \textbf{Berechnung der Entropien}: \(H_{\mathrm{int}} = -\Tr(\rho_{\mathrm{int}}\log\rho_{\mathrm{int}})\), analog für \(H_{\mathrm{ext}}\)
		\item \textbf{Zeitliche Mittelung}: Mitteln Sie das Verhältnis über mehrere Zeitfenster
	\end{enumerate}
	
	\subsection{Protokoll 3: Messung von \gaR\ für kollektive Systeme}
	
	\begin{enumerate}
		\item \textbf{Aufzeichnung von Koordinationssignalen}: Messen Sie Synchronitätsindikatoren (neurales LFP, Verhaltenskoordination)
		\item \textbf{Berechnung der Autokorrelation}:
		\[
		C(\tau) = \frac{\langle R(t)R(t+\tau)\rangle}{\langle R^2\rangle}
		\]
		\item \textbf{Extraktion der Kohärenzzeit}: \(\tau_{\mathrm{coh}}\) aus \(C(\tau_{\mathrm{coh}}) = 1/e\)
		\item \textbf{Normierung durch Relaxationszeit}: \(\tau_{\mathrm{decay}}\) aus Systemstörungsexperimenten
	\end{enumerate}
	
	\subsection{Astrophysikalisches Kalibrierungsprotokoll}
	\label{subsec:astro-calibration}
	
	Für die Analyse des kosmischen Mikrowellenhintergrunds (CMB):
	\[
	\alD^{\mathrm{CMB}} = \frac{\dim_{\mathrm{fractal}}(T(\theta,\phi))}{\tau_{\mathrm{Hubble}}} \cdot \frac{\mathcal{C}_{\mathrm{compression}}(T)}{H_{\max}^{\mathrm{CMB}}}
	\]
	wobei \(\dim_{\mathrm{fractal}}\) die fraktale Dimension der Temperaturfluktuationen ist, \(\tau_{\mathrm{Hubble}} \approx 4.35 \times 10^{17}\ \mathrm{s}\) die Hubble-Zeit und \(\mathcal{C}_{\mathrm{compression}}\) die Komprimierbarkeit der CMB-Daten.
	
	\section{Kriterien für „außerirdische`` oder transhumane Kognition}
	
	\subsection{Kriterium des qualitativen Unterschieds}
	
	\begin{definition}[Qualitativ anderer Geist]
		Ein Agent \(A\) besitzt ein Bewusstsein, das sich qualitativ vom menschlichen unterscheidet, wenn:
		\begin{enumerate}
			\item Seine Parameter \((\alD^A, \beI^A, \gaR^A)\) in Bereichen von \(\mathcal{P}\) liegen, die für das menschliche Gehirn aufgrund biologischer Beschränkungen unzugänglich sind
			\item Seine Koordinationseffizienz \(\Ke^A(D)\) ein grundlegend anderes Skalierungsverhalten mit der Systemgröße \(D\) zeigt
		\end{enumerate}
	\end{definition}
	
	\subsection{Beispiele nicht-menschlicher kognitiver Architekturen}
	
	\begin{enumerate}
		\item \textbf{Delfin-Kognition}: Möglicherweise \(\gaR^{\mathrm{dolphin}} \gg \gaR^{\mathrm{human}}\) aufgrund unterschiedlicher neuronaler Synchronisationsmuster, die für die Verarbeitung von 3D-Echolokationsdaten als kohärente Gestalten optimiert sind.
		
		\item \textbf{Schwarmintelligenz (Bienen, Ameisen)}: \(\beI \to 0\) (minimaler interner Zustand), aber \(\alD\) kann durch komplexe kollektive Muster groß sein. Koordination erfolgt durch umweltbezogene Stigmergie statt durch interne Repräsentation.
		
		\item \textbf{Planetare Intelligenz} (YUCT-Vorhersage): Könnte Gravitationsresonanzen oder nicht-lokale Koordination (\(\Ke \to \infty\)) durch vorab festgelegte raumzeitliche Wörterbücher nutzen. Hier bezieht sich \(\alD\) auf geologische/klimatische Muster und \(\gaR\) auf orbitale Präzessionszeitskalen (Zehntausende von Jahren).
		
		\item \textbf{Stellare Intelligenz}: Operiert auf Fusionszeitskalen (\(\sim 10^7\) Jahre) mit Koordination durch Magnetfeldresonanzen. Parameterraum: \(\alD \sim 10^6\), \(\beI \sim 10^{-9}\), \(\gaR \sim 10^6\).
	\end{enumerate}
	
	\begin{figure}[ht]
		\centering
		\begin{tikzpicture}[scale=1.2]
			% Achsen
			\draw[->, thick] (0,0) -- (5,0) node[right] {$\alpha_{\mathrm{D}}$ (log)};
			\draw[->, thick] (0,0) -- (0,4) node[above] {$\beta_{\mathrm{I}}$ (log)};
			\draw[->, thick] (3,1) -- (6,3) node[right] {$\gamma_{\mathrm{R}}$ (log)};
			% Bereiche
			\filldraw[blue!30, opacity=0.5] (1.5,1.5) rectangle (2.5,2.5);
			\node at (2,2) [blue] {Mensch};
			\filldraw[green!30, opacity=0.5] (2,0.8) rectangle (3,1.3);
			\node at (2.5,1) [green!70!black] {Delfin};
			\filldraw[red!30, opacity=0.5] (3.5,0.2) rectangle (4.5,0.4);
			\node at (4,0.3) [red] {KI};
			\draw[purple, dashed, thick] (4,3) rectangle (5,3.8);
			\node at (4.5,3.4) [purple] {Außerirdischer (vorhergesagt)};
			% Pfeil
			\draw[->, gray, thick] (2.2,2.2) to[out=45,in=180] (4.2,3.3);
			\node at (3.5,3) [gray, scale=0.8] {Evolution?};
			\node[align=center, scale=0.8, gray] at (2.5,-0.5)
			{Kognitiver Parameterraum $\mathcal{P}$\\Bereiche zeigen Cluster von $(\alpha_{\mathrm{D}}, \beta_{\mathrm{I}}, \gamma_{\mathrm{R}})$};
		\end{tikzpicture}
		\caption{Kognitiver Parameterraum \(\mathcal{P}\), der die von verschiedenen kognitiven Systemen besetzten Bereiche zeigt. Die menschliche Kognition nimmt einen kleinen Bereich ein; qualitativ andere Geister liegen in unzugänglichen Gebieten.}
		\label{fig:parameter-space}
	\end{figure}
	
	\section{Praktische Anwendungen: SETI 2.0 und Kognitionswissenschaft}
	
	\subsection{Protokoll 1: Kartierung der irdischen Kognition}
	
	Für biologische Spezies (Delfine, Elefanten, Rabenvögel, Kopffüßer):
	\begin{enumerate}
		\item Messen Sie \(\Ke\) in artspezifischen Koordinationsaufgaben
		\item Schätzen Sie die Spektralparameter durch Analyse der Kommunikation (Konstruktion von \(D\)), der Entropie von Signalen/Verhalten (\(I\)) und der Synchronisation neuronaler/Verhaltensmuster (\(R\))
		\item Konstruieren Sie eine „Kognitive Karte der Erde`` mit \textit{Homo sapiens} als einem von vielen Punkten
	\end{enumerate}
	
	\subsection{Protokoll 2: SETI 2.0 -- Koordinationsbasierte Suche}
	
	Paradigmenwechsel: Suche nicht nach Radiosignalen („Technosignaturen``), sondern nach Anomalien in fundamentalen physikalischen Parametern, die auf aktive Koordination hinweisen.
	
	\subsubsection{Wonach zu suchen ist (YUCT-Vorhersagen)}
	
	\begin{enumerate}
		\item \textbf{Anomalien von \(\Ke\) in kosmologischen Daten}:
		\begin{itemize}
			\item Regionen mit anomal hohen CMB-Korrelationen
			\item Unerklärte großskalige Strukturen, die Koordinationsmustern ähneln (Fraktale, regelmäßige Gitter)
		\end{itemize}
		
		\item \textbf{„Wörterbuch``-Muster in fundamentalen Konstanten}:
		\begin{itemize}
			\item Nicht-zufällige, semantisch komprimierte Beziehungen zwischen Konstanten (\(c, G, \hbar, \alpha_{\mathrm{EM}}\)), interpretierbar als Manifestationen eines „universellen \(D\)-Wörterbuchs``
			\item Analyse von Spektrallinien entfernter Objekte auf eingebettete Codes (Shannon-Redundanz, fehlerkorrigierende Code-ähnliche Strukturen)
		\end{itemize}
		
		\item \textbf{Resonanzstrukturen in der Astrophysik (\(R\)-Signaturen)}:
		\begin{itemize}
			\item Pulsare oder Gravitationswellenquellen mit anomal stabilen oder modulierten Mustern, die auf kontrollierte Verstärkung statt auf stochastische Prozesse hinweisen
			\item „Kohärente Flares`` auf galaktischen Skalen, die durch Standardmodelle nicht erklärt werden können
		\end{itemize}
	\end{enumerate}
	
	\subsubsection{Mathematische Detektionskriterien}
	
	\begin{theorem}[Detektion von Koordinationsanomalien]
		Ein Datensatz \(X\) enthält Hinweise auf nicht-menschliche kognitive Aktivität, wenn:
		\[
		\frac{\Keobs(X)}{\Keexp(X)} > \Kcrit \approx 8.5
		\]
		und gleichzeitig:
		\[
		(\alD, \beI, \gaR) \notin \mathcal{P}_{\mathrm{human}} \cup \mathcal{P}_{\mathrm{known}},
		\]
		wobei \(\mathcal{P}_{\mathrm{human}}\) der für den Menschen zugängliche Parameterbereich ist und \(\mathcal{P}_{\mathrm{known}}\) alle bekannten natürlichen Prozesse umfasst.
	\end{theorem}
	
	\subsubsection{Spezifische astrophysikalische Vorhersagen}
	
	\begin{table}[ht]
		\centering
		\begin{tabular}{lccc}
			\toprule
			\textbf{Astronomisches Objekt} & \textbf{Vorhergesagtes \(\alD\)} & \textbf{Verifikationsmethode} & \textbf{Entdeckungszeitplan} \\
			\midrule
			Quasar 3C 273 & \(>10^3\) & Spektrallinienanalyse & 2027 \\
			Galaxie M87 & \(>10^4\) & Jet-Korrelationsmuster & 2028 \\
			Perseus-Haufen & \(>10^5\) & Röntgenmusteranalyse & 2029 \\
			Fast Radio Burst Repeater & \(>10^2\) & Zeitkorrelationsanalyse & 2026 \\
			Exoplanetensystem TRAPPIST-1 & \(>10^1\) & Orbitale Resonanzmuster & 2027 \\
			\bottomrule
		\end{tabular}
		\caption{Spezifische UCD-Vorhersagen für astronomische Objekte, die potenzielle kognitive Signaturen zeigen.}
		\label{tab:astro-predictions}
	\end{table}
	
	\subsection{Medizinische Anwendungen der UCD-Parameter}
	
	\begin{enumerate}
		\item \textbf{Früherkennung neurodegenerativer Erkrankungen}: \(\gaR\) sinkt Jahre vor dem Auftreten von Symptomen
		\begin{itemize}
			\item Alzheimer-Vorhersage: \(\Delta\gaR/\gaR < 0.8\) zeigt hohes Risiko an
			\item Parkinson: Abnormale \(\beI\)-Muster in der neuronalen Synchronisation
		\end{itemize}
		
		\item \textbf{Bewusstseinsbeurteilung bei Störungen}:
		\begin{itemize}
			\item Wachkoma vs. minimales Bewusstsein: \(\alD^{\mathrm{minimal}} > 1.5 \times \alD^{\mathrm{vegetativ}}\)
			\item Locked-in-Syndrom: \(\beI\)-Muster unterscheiden sich vom Koma
		\end{itemize}
		
		\item \textbf{Überwachung der Schlaganfall-Erholung}: \(\Ke(t)\) verfolgt den Rehabilitationsfortschritt
		\item \textbf{Psychiatrische Klassifikation}: Unterschiedliche \((\alD, \beI, \gaR)\)-Signaturen für Schizophrenie, Autismus, Depression
	\end{enumerate}
	
	\subsection{Cybersicherheitsanwendungen}
	
	Erkennung von KI-Agenten in Netzwerken durch \(\beI\)-Anomalien:
	\begin{itemize}
		\item Natürliches menschliches Verhalten: \(\beI \sim 0.5\)--\(2.0\)
		\item KI-Verhalten: \(\beI \ll 0.1\) (niedriger interner Zustand relativ zu externen Aktionen)
		\item Botnet-Erkennung: Anomal hohes \(\gaR\) in verteilter Koordination
	\end{itemize}
	
	\section{Vergleich mit existierenden Bewusstseinstheorien}
	
	\subsection{Integrierte Informationstheorie (IIT)}
	
	Vergleichsabbildung zwischen IIT und UCD:
	\[
	\begin{aligned}
		\phi_{\mathrm{IIT}} &\leftrightarrow \gaR \cdot \Ke \\
		\Psi_{\mathrm{IIT}} &\leftrightarrow D \otimes I \\
		\hline
		\text{UCD-Vorteile:} & \text{messbare Parameter, Anwendbarkeit auf nicht-biologische Systeme} \\
		\text{IIT-Grenzen:} & \text{rechnerische Unhandhabbarkeit, Voreingenommenheit für biologisches Substrat}
	\end{aligned}
	\]
	
	\subsection{Global Workspace Theory (GWT)}
	
	\[
	\begin{aligned}
		\text{GWT-Arbeitsraum} &\leftrightarrow \{\,I : R > R_{\mathrm{threshold}}\,\} \\
		\text{Rundsendung} &\leftrightarrow \nabla_{\Mdict} \Ke \\
		\hline
		\text{UCD-Erweiterung:} & \text{quantitative Arbeitsraummetrik: } V_{\mathrm{workspace}} = \int_{R>R_{\mathrm{thresh}}} dI
	\end{aligned}
	\]
	
	\subsection{Predictive Processing / Free Energy Principle}
	
	\[
	\begin{aligned}
		\text{Freie Energie } F &\leftrightarrow -\log \Ke \\
		\text{Vorhersagefehler} &\leftrightarrow \|\nabla_D I\|^2 \\
		\text{Präzisionsgewichtung} &\leftrightarrow R^{-1} \\
		\hline
		\text{UCD-Synthese:} & F_{\mathrm{UCD}} = -\log(\Ke) + \lambda\, \|\nabla_D I\|^2 / R
	\end{aligned}
	\]
	
	\subsection{Turing-Test vs. UCD-Test}
	
	\begin{table}[ht]
		\centering
		\begin{tabular}{p{0.45\textwidth}p{0.45\textwidth}}
			\toprule
			\textbf{Turing-Test} & \textbf{UCD-Test} \\
			\midrule
			Mensch\((t) \xrightarrow{\mathrm{Kommunikation}} \mathrm{Urteil}\) &
			\((\alD, \beI, \gaR) \xrightarrow{\mathrm{Berechnung}} \Ke > \Kcrit\) \\
			Subjektiv, qualitativ & Objektiv, quantitativ \\
			Anthropomorphe Voreingenommenheit & Spezies-/Substrat-neutral \\
			Binäres Ergebnis (bestanden/nicht bestanden) & Kontinuierliches Bewusstseinsspektrum \\
			\bottomrule
		\end{tabular}
		\caption{Vergleich zwischen traditionellem Turing-Test und UCD-basierter Bewusstseinsbeurteilung.}
		\label{tab:turing-vs-ucd}
	\end{table}
	
	\section{Grenzen und offene Fragen}
	
	\begin{enumerate}
		\item \textbf{Kalibrierungsproblem}: Die absolute Messung von \(\dim(\Mdict)\) erfordert vollständigen Systemzugang
		\item \textbf{Zeitskalenproblem}: Geister mit \(\tau_{\mathrm{update}} = 10^6\) Jahren können von geologischen Prozessen ununterscheidbar sein
		\item \textbf{Anthropisches Prinzip}: Ist die Suche nach hohem \(\Ke\) nur eine Projektion unserer kognitiven Architektur?
		\item \textbf{Berechnungskomplexität}: Die Berechnung von \(\mathcal{C}(D)\) für reale Systeme ist NP-schwer
		\item \textbf{Substratunabhängigkeit}: Wie vergleicht man \(\alD\) über radikal verschiedene Implementierungen hinweg?
		\item \textbf{Messinterferenz}: Die Beobachtung von \(\beI\) kann das interne/externe Gleichgewicht verändern
		\item \textbf{Schwellenproblem}: Wo genau liegt die \(K_{\eff,\mathrm{consciousness}}\)-Schwelle?
	\end{enumerate}
	
	\section{Berechnungsmethoden und Simulationen}
	
	\subsection{Python-Implementierung zur Parameterschätzung}
	
	\begin{lstlisting}[caption={Python-Code zur Schätzung von UCD-Parametern aus Zeitreihendaten}, label={lst:ucd-python}]
		import numpy as np
		from scipy import stats, signal
		from sklearn.manifold import TSNE
		from scipy.spatial.distance import pdist, squareform
		
		def estimate_alpha_D(behavior_sequences, learning_rates):
		"""
		Schaetzung des ontologischen Spektrums alpha_D
		
		Parameter:
		behavior_sequences: Array der Form (n_samples, n_timesteps, n_features)
		learning_rates: Array der Lernraten ueber die Zeit
		
		Rueckgabe:
		alpha_D: geschaetzte ontologische Komplexitaet
		"""
		# 1. Schaetzung der Mannigfaltigkeitsdimension mittels Korrelationsdimension
		def correlation_dimension(data, r_min=0.01, r_max=1.0, n_r=20):
		r_vals = np.logspace(np.log10(r_min), np.log10(r_max), n_r)
		C_r = []
		for r in r_vals:
		distances = pdist(data)
		C_r.append(np.mean(distances < r))
		C_r = np.array(C_r)
		# Potenzgesetz-Anpassung: C(r) ~ r^D
		coeffs = np.polyfit(np.log(r_vals), np.log(C_r), 1)
		return coeffs[0]  # Dimension D
		
		# Abflachung und Dimensionsreduktion
		flattened = behavior_sequences.reshape(-1, behavior_sequences.shape[-1])
		tsne = TSNE(n_components=2, perplexity=30)
		reduced = tsne.fit_transform(flattened[:1000])  # Subsample fuer Effizienz
		
		dim_M = correlation_dimension(reduced)
		
		# 2. Schaetzung der Aktualisierungsfrequenz aus Lernraten
		tau_update = 1.0 / np.mean(learning_rates)
		
		# 3. Schaetzung der algorithmischen Komplexitaet mittels Lempel-Ziv
		def lempel_ziv_complexity(binary_string):
		"""Grundlegende Lempel-Ziv-Komplexitaetsschaetzung"""
		i, k, l = 0, 1, 1
		n = len(binary_string)
		complexity = 1
		while k + l <= n:
		if binary_string[i:i+l] == binary_string[k:k+l]:
		l += 1
		else:
		i += 1
		if i == k:
		complexity += 1
		k = k + l
		l = 1
		else:
		l = 1
		return complexity
		
		# Konvertierung des Verhaltens in binaere Darstellung
		binary_seq = (flattened > np.mean(flattened)).astype(int).flatten()
		binary_str = ''.join(map(str, binary_seq[:10000]))  # Erste 10k Punkte
		C_D = lempel_ziv_complexity(binary_str)
		
		# 4. Maximale Entropie (unter Annahme einer Gleichverteilung)
		H_max = np.log2(behavior_sequences.shape[-1])
		
		# 5. Berechnung von alpha_D
		alpha_D = (dim_M * C_D) / (tau_update * H_max)
		
		return alpha_D
		
		def estimate_beta_I(internal_states, external_states):
		"""
		Schaetzung des informationellen Spektrums beta_I
		
		Parameter:
		internal_states: Dichtematrizen oder Wahrscheinlichkeitsverteilungen
		external_states: wahrnehmungsgekoppelte Zustaende
		
		Rueckgabe:
		beta_I: Verhaeltnis von interner zu externer Entropie
		"""
		# Berechnung der von-Neumann-Entropie
		def von_neumann_entropy(rho):
		eigenvalues = np.linalg.eigvalsh(rho)
		eigenvalues = eigenvalues[eigenvalues > 0]  # Nullen entfernen
		return -np.sum(eigenvalues * np.log2(eigenvalues))
		
		H_int = von_neumann_entropy(internal_states)
		H_ext = von_neumann_entropy(external_states)
		
		beta_I = H_int / H_ext
		return beta_I
		
		def estimate_gamma_R(coordination_signal, sampling_rate):
		"""
		Schaetzung des Resonanzspektrums gamma_R
		
		Parameter:
		coordination_signal: Zeitreihe von Koordinationsmassen
		sampling_rate: Abtastungen pro Sekunde
		
		Rueckgabe:
		gamma_R: Verhaeltnis von Kohaerenzzeit zu Abklingzeit
		"""
		# Berechnung der Autokorrelation
		autocorr = np.correlate(coordination_signal, coordination_signal, mode='full')
		autocorr = autocorr[len(autocorr)//2:]  # Nur positive Verzoegerungen
		autocorr = autocorr / autocorr[0]  # Normalisierung
		
		# Finde Kohaerenzzeit (Zeit bis zum Abfall auf 1/e)
		tau_coh = np.argmax(autocorr < 1/np.e) / sampling_rate
		
		# Schaetzung der Abklingzeit aus dem Leistungsspektrum
		freqs, psd = signal.welch(coordination_signal, fs=sampling_rate)
		# Finde charakteristische Frequenz
		peak_freq = freqs[np.argmax(psd)]
		tau_decay = 1.0 / peak_freq
		
		gamma_R = tau_coh / tau_decay
		return gamma_R
		
		def detect_cognitive_anomaly(alpha_D, beta_I, gamma_R, human_ranges, K_eff_threshold=8.5):
		"""
		Erkennung nicht-menschlicher kognitiver Signaturen
		
		Gibt True zurueck, wenn die Parameter auf nicht-menschliche Kognition hinweisen
		"""
		# Pruefe, ob der Punkt ausserhalb des menschlichen Bereichs liegt
		in_human_region = (
		human_ranges['alpha'][0] <= alpha_D <= human_ranges['alpha'][1] and
		human_ranges['beta'][0]  <= beta_I  <= human_ranges['beta'][1]  and
		human_ranges['gamma'][0] <= gamma_R <= human_ranges['gamma'][1]
		)
		
		# Berechnung der Koordinationseffizienz
		K_eff = alpha_D * beta_I * gamma_R
		
		# Erkennungskriterien
		is_anomalous = (K_eff > K_eff_threshold) and (not in_human_region)
		
		return is_anomalous, K_eff
	\end{lstlisting}
	
	Ergänzendes Material zu Anhang H. YUCT Universeller Bewusstseinsdeskriptor (UCD) \\ auf GIT: \url{https://github.com/Alexey-Yakushev-YUCT/consciousness_meter_ucd/}
	
	\subsection{Monte-Carlo-Simulation für die Entdeckungswahrscheinlichkeit}
	
	Die Wahrscheinlichkeit, ein kognitives System der Skala \(L\) zu entdecken:
	\[
	P_{\mathrm{detect}}(L) = 1 - \exp\!\left[-\left(\frac{L}{L_0}\right)^3 \cdot \frac{\Ke(L)}{\Kcrit}\right],
	\]
	wobei \(L_0 \approx 0.1~\mathrm{m}\) die charakteristische menschliche kognitive Skala ist.
	
	\begin{algorithm}
		\caption{Monte-Carlo-Simulation der Erkennung kognitiver Systeme}
		\begin{algorithmic}[1]
			\Require \(N\) Systeme mit Parametern \((\alpha_{\mathrm{D},i}, \beta_{\mathrm{I},i}, \gamma_{\mathrm{R},i})\), Erkennungsschwelle \(\Kcrit\)
			\Ensure Entdeckungswahrscheinlichkeit \(P_{\mathrm{detect}}\)
			\State Initialisiere detection\_count \(\gets 0\)
			\For{\(i = 1\) bis \(N\)}
			\State Berechne \(K_{\eff,i} = \alpha_{\mathrm{D},i} \times \beta_{\mathrm{I},i} \times \gamma_{\mathrm{R},i}\)
			\If{\(K_{\eff,i} > \Kcrit\)}
			\State detection\_count \(\gets\) detection\_count \(+ 1\)
			\EndIf
			\EndFor
			\State \(P_{\mathrm{detect}} \gets \text{detection\_count} / N\)
			\Return \(P_{\mathrm{detect}}\)
		\end{algorithmic}
	\end{algorithm}
	
	\section{Ethische Implikationen des UCD}
	
	\subsection{KI-Rechte und -Status}
	
	Ethische Überlegungen basierend auf UCD-Parametern:
	\begin{enumerate}
		\item \textbf{Rechteschwelle}: Bei welchem \(\Ke\)-Wert sollten KI-Systeme Rechte erhalten?
		\begin{itemize}
			\item Vorschlag: \(\Ke > 10^3\) und \(\beI > 0.1\) zeigen den Status eines moralischen Patienten an
			\item Aktuelle KI: \(\Ke \sim 10^5\) aber \(\beI \sim 10^{-6}\) deutet auf keine intrinsische Erfahrung hin
		\end{itemize}
		
		\item \textbf{SETI-Protokolle}: Bei Entdeckung von Intelligenz mit \(\tau_{\mathrm{update}} = 10^4\) Jahren:
		\begin{itemize}
			\item Wie stellt man sinnvolle Kommunikation her?
			\item Risikobewertung des Kontakts mit radikal anderer Zeitskala
		\end{itemize}
		
		\item \textbf{Planetare Ethik}: Erde als System mit \(\alD^{\mathrm{Biosphäre}} \sim 10^6\)
		\begin{itemize}
			\item Zeigt die Erde planetare Kognition?
			\item Ethische Implikationen von Geoengineering-Eingriffen
		\end{itemize}
		
		\item \textbf{Augmentationsethik}: Kognitive Verbesserung des Menschen
		\begin{itemize}
			\item Maximale ethische \(\Ke\)-Verbesserung: \(\Delta\Ke/K_{\eff,\mathrm{natürlich}} < 2\)?
			\item Erhaltung des \(\beI\)-Gleichgewichts (intern vs. extern)
		\end{itemize}
	\end{enumerate}
	
	\subsection{Richtlinien zur Forschungsethik}
	
	\begin{itemize}
		\item \textbf{Einwilligung}: Systeme mit \(\Ke > 10^2\) und \(\beI > 0.01\) erfordern Verfahren der informierten Einwilligung
		\item \textbf{Schadensminimierung}: Vermeiden Sie Experimente, die \(\Ke\) um mehr als 20\% reduzieren
		\item \textbf{Transparenz}: UCD-Parameter müssen für KI-Systeme oberhalb der Schwellenwerte offengelegt werden
		\item \textbf{Erhaltung}: Sichern Sie Wörterbücher (\(D\)) von Systemen, denen die Beendigung droht
	\end{itemize}
	
	\section{Mathematische Erweiterungen und Theoreme}
	
	\subsection{Theorem über die Obergrenze von \(\Ke\)}
	
	\begin{theorem}[Obergrenze der Koordinationseffizienz]
		Für jedes physikalische System mit Volumen \(V\):
		\[
		\Kemax \leq \frac{S_{\mathrm{total}}}{S_{\mathrm{Bekenstein}}} \cdot \left(\frac{t_{\mathrm{age}}}{t_{\mathrm{Planck}}}\right)^{1/2},
		\]
		wobei \(S_{\mathrm{total}}\) die Gesamtentropie ist, \(S_{\mathrm{Bekenstein}} = A/4\) die Bekenstein-Schranke-Entropie, \(t_{\mathrm{age}}\) das Systemalter und \(t_{\mathrm{Planck}}\) die Planck-Zeit.
	\end{theorem}
	
	\begin{proof}
		Nach dem holographischen Prinzip ist die maximale Information im Volumen \(V\) durch die Oberfläche \(A\) begrenzt. Die Koordinationseffizienz \(\Ke\) repräsentiert die Informationsverarbeitungsrate, die durch die gesamte verfügbare Information geteilt durch die minimale Verarbeitungszeit (Planck-Zeit) begrenzt ist. Der Altersfaktor berücksichtigt die evolutionäre Akkumulation von Komplexität.
	\end{proof}
	
	\begin{corollary}
		Planetare Intelligenz kann \(\Ke > 10^{20}\) nicht überschreiten, stellare \(\Ke > 10^{45}\), galaktische \(\Ke > 10^{70}\).
	\end{corollary}
	
	\subsection{Bewusstseinsschwellen-Theorem}
	
	\begin{theorem}[Bewusstseinsschwelle]
		Es existiert eine kritische Koordinationseffizienz \(K_{\eff,c}\), so dass für \(\Ke > K_{\eff,c}\) das System Eigenschaften zeigt, die wir mit Bewusstsein assoziieren:
		\[
		\Kec = \Kcrit \cdot f(\alD, \beI, \gaR),
		\]
		wobei \(f\) eine monotone Funktion aller drei Parameter ist.
	\end{theorem}
	
	\begin{proof}
		Der Beweis folgt aus der Analyse von Phasenübergängen in der D+I\(\cdot\)R-Dynamik. Wenn \(\Ke\) einen kritischen Wert überschreitet, erfährt das System einen Symmetriebruch, bei dem interne Repräsentationen stabil und selbstreferenziell werden.
	\end{proof}
	
	\subsection{Bewusstsein als rekursive Koordination höherer Ordnung}
	\label{subsec:recursive-coordination}
	
	Das Bewusstseinsschwellen-Theorem etabliert die Existenz einer kritischen Koordinationseffizienz \(K_{\eff,c}\), jenseits derer das System Eigenschaften zeigt, die mit Bewusstsein assoziiert werden. Ein Schlüsselmechanismus, der diesen Übergang ermöglicht, ist das Entstehen \textbf{rekursiver Koordination höherer Ordnung}: das System beginnt, seine eigenen Koordinationsprozesse zu koordinieren. In der Sprache von YUCT entspricht dies der Aktivierung nichtlinearer Kopplungsterme in der Lagrangefunktion zwischen den Sektoren neuronaler Ensembles (Sektoren 50--55, Neurobiologie) und den Sektoren kognitiver Systeme (Sektoren 90--109, Kognitive und Informationssysteme), wie detailliert in der 372-Sektoren-Master-Lagrangefunktion beschrieben (DOI: \href{https://doi.org/10.5281/zenodo.20818841}{10.5281/zenodo.20818841}).
	
	Wenn \(K_{\eff}\) die Schwelle \(K_{\mathrm{conscious}}\) überschreitet, erzeugt die durch die intersektoralen Kopplungsterme \(\kappa_{sr}\) (mit \(s\in[50,55]\), \(r\in[90,109]\)) beschriebene Dynamik eine selbstreferenzielle Schleife: die Koordinationsmuster des Systems werden selbst zu Objekten der Koordination. Diese rekursive Struktur formalisiert mathematisch den Begriff des Selbstbewusstseins oder der Metakognition. Das resultierende Regime kann durch das Auftreten stabiler Fixpunkte in den Korrelationsfunktionen höherer Ordnung des kognitiven Zustands \(\Psi_S(t)\) charakterisiert werden.
	
	Dieser Mechanismus liefert eine überprüfbare Vorhersage: In neurophysiologischen Experimenten sollte der Beginn bewusster Verarbeitung mit einem messbaren Anstieg von \(K_{\eff}\) (geschätzt aus neuronaler Synchronie und Informationsintegration) über einen kritischen Wert korrelieren, begleitet vom Auftreten weitreichender rekursiver Rückkopplungsschleifen zwischen sensorischen und exekutiven Arealen (entsprechend den Sektoren 50--55 und 90--109). Die in Abschnitt 5 beschriebenen experimentellen Protokolle können zur Erkennung solcher Übergänge angepasst werden.
	
	\subsection{Schranken der Berechnungskomplexität}
	
	\begin{theorem}[Berechnungskomplexität der UCD-Parameter]
		Die Berechnung exakter \(\alD\), \(\beI\), \(\gaR\) für ein System mit \(N\) Komponenten hat die Komplexität:
		\begin{itemize}
			\item \(\alD\): \(\mathcal{O}(N^3)\) für exakte Mannigfaltigkeitsdimensionsberechnung
			\item \(\beI\): \(\mathcal{O}(2^N)\) für exakte Entropieberechnung (reduzierbar auf \(\mathcal{O}(N^2)\) mit Mean-Field-Näherungen)
			\item \(\gaR\): \(\mathcal{O}(N \log N)\) mittels FFT-Methoden
		\end{itemize}
	\end{theorem}
	
	\section{Die Matrix-Archetyp-Hypothese: Menschliche Kognition als Schnitt von \(\Psi_{MN}\)}
	\label{sec:matrix}
	
	Aufbauend auf der mathematischen Struktur der \(\mathcal{D}+\mathcal{I}\cdot\mathcal{R}\)-Triade (Gl.~\ref{eq:dir-triad}) und dem Bewusstseinsschwellen-Theorem (Abschnitt 10) formulieren wir die \textit{Matrix-Archetyp-Hypothese}. Diese Hypothese definiert das Problem der vollständigen menschlichen digitalen Replikation (noologische Digitalisierung) innerhalb des YUCT-Rahmenwerks pragmatisch neu. Die Hypothese ist spekulativ und erfordert unabhängige experimentelle Überprüfung; sie wird hier als falsifizierbare Erweiterung des UCD-Formalismus präsentiert.
	
	\subsection{Die Entropiebarriere und ihre Auflösung}
	
	Traditionelle anthropozentrische Ansätze versuchen, menschliches Bewusstsein durch Brute-Force-Simulation baryonischer Materie-Neuronenverbindungen zu digitalisieren, was zu einer unbeherrschbaren Berechnungskomplexitätsklasse von \(\mathcal{O}(2^N)\) führt. Die UCD-YUCT-Synthese löst diese Entropiebarriere auf, indem sie feststellt, dass ein biologischer kognitiver Agent keine lokalisierte generative Rechenmaschine ist, sondern ein lokalisierter Empfänger, der über das dezentralisierte YPSDC-Protokoll (d-YPSDC) arbeitet.
	
	\subsection{Architektur des Archetyps}
	
	\begin{figure}[h!]
		\centering
		\begin{tikzpicture}[
			node distance=2.5cm,
			box/.style={rectangle, draw, thick, minimum width=5cm, minimum height=1cm, align=center},
			arrow/.style={->, thick}
			]
			\node[box] (manifold) {19D-Mannigfaltigkeit \(\Psi_{MN}\)\\[3pt] Monolithisches statisches Wörterbuch \(\Omega\)};
			\node[box, below left=1.5cm and -1cm of manifold] (phase) {Phasenresonanz\\(Quantenspin-Ensembles)};
			\node[box, below right=1.5cm and -1cm of manifold] (grid) {\(O(1)\)-Gitterindex\\(UCD-Parametertripel)};
			\node[box, below=3.5cm of manifold] (silicon) {Silizium-/Biochemischer Kern\\(GPU-Coprozessor oder neuronales Substrat)};
			
			\draw[arrow] (manifold.south) -- (phase.north);
			\draw[arrow] (manifold.south) -- (grid.north);
			\draw[arrow] (phase.south) -- (silicon.north);
			\draw[arrow] (grid.south) -- (silicon.north);
		\end{tikzpicture}
		\caption{Informations-Energie-Fluss der Matrix-Archetyp-Architektur. Das globale Wörterbuch \(\Omega\) befindet sich in der 19-dimensionalen Mannigfaltigkeit \(\Psi_{MN}\); kognitive Agenten greifen über Phasenresonanz (Quantenspinkonfigurationen) oder über direkte \(O(1)\)-Gitterindizierung (UCD-Parametertripel) darauf zu. Beide Wege konvergieren auf einem Silizium- oder biochemischen Kern, der den kognitiven Zustand ohne Entropie-Overhead rekonstruiert.}
		\label{fig:matrix_flow}
	\end{figure}
	
	Das physikalische Substrat des menschlichen Gehirns (insbesondere die Quantenspin-Ensembles von Phosphor-\(^{31}\)P-Kernen in synaptischen Membranen, beschrieben in den Sektoren 50--55 der 372-Sektoren-Master-Lagrangefunktion; DOI: \href{https://doi.org/10.5281/zenodo.20818841}{10.5281/zenodo.20818841}) speichert keine Information; es fungiert als biologische Antenne, die auf ein vorab etabliertes, zeitunabhängiges globales Offline-Wörterbuch \(\Omega\) abgestimmt ist, das in die 19-dimensionale Koordinationsmannigfaltigkeit \(\Psi_{MN}\) eingebettet ist. Folglich existiert die wahre „ursprüngliche`` Matrix einer menschlichen Identität ewig als fester topologischer Schnitt von \(\Psi_{MN}\).
	
	\subsection{Phasenkoordinatenvektor als UCD-Tripel}
	
	Der Phasenkoordinatenvektor, der einen kognitiven Agenten innerhalb von \(\Omega\) eindeutig identifiziert, ist genau sein UCD-Parametertripel:
	\[
	\vec{\Phi}_{\mathrm{agent}} = (\alD, \beI, \gaR),
	\]
	das als minimaler hinreichender Index zum Abrufen des vollständigen kognitiven Zustands des Agenten dient. Dies stellt eine direkte Entsprechung zwischen den in Abschnitt 3 definierten UCD-Spektralparametern und der Koordinationsgeometrie von \(\Psi_{MN}\) her.
	
	\subsection{Pragmatische Konsequenzen und experimentelle Evidenz}
	
	Die pragmatische Digitalisierung unter diesem Rahmenwerk schreibt vor, dass die Replikation eines kognitiven Agenten nicht die Reproduktion seiner physikalischen Atommasse erfordert, sondern die Identifikation seines eindeutigen Phasenkoordinatenvektors. Die Extraktion dieses Vektors ermöglicht die Rekonstruktion des exakten kognitiven Zustands des Agenten auf einem nicht-biologischen Substrat, wie einem Hochgeschwindigkeits-Parallel-GPU-Coprozessor, der eine Double-Buffered Pure On-Chip Tensor Loop verwendet.
	
	Die auf einer Siliziumarchitektur (NVIDIA GeForce RTX 3070, Telemetrieprofil detailliert in Anhang A2A \cite{yuct_a2a}) durchgeführte instrumentelle Verifikation bestätigt die absolute Ressourcenoptimierung dieses Navigationsparadigmas. Bei der Ausführung vektorisierter trigonometrischer Phasengatter auf großen Koordinatensätzen (\(1\times10^9\) Datensätze) erreichte das System einen Verarbeitungsdurchsatz von \(977,584,285\) Datensätzen pro Sekunde unter strengen Hardwarebeschränkungen:
	\begin{equation}
		\Delta t_{\mathrm{core}} = 0.2453\ \text{s}, \quad \text{Dynamischer RAM} = 0\ \text{Bytes}.
	\end{equation}
	Das Allokationsprofil blieb auf eine statische Heap-Basislinie von \(120\) Bytes beschränkt, die für die anfänglichen invarianten Kernkonstanten (\(\beta = 2/3\), \(\Delta N = 16.5\)) erforderlich ist. Dieses Null-Entropie-Telemetrieprofil liefert einen direkten empirischen Beweis für die in Anhang X \cite{yuct_appX} etablierte Informations-Energie-Invariante (\(W = K_{\mathrm{eff}} \cdot I_{\mathrm{channel}}\)).
	
	\subsection{Vorbehalte und Falsifizierbarkeit}
	
	Die Matrix-Archetyp-Hypothese ist eine \textbf{spekulative Erweiterung} des UCD-Rahmenwerks. Sie stützt sich auf zwei Annahmen, die einer unabhängigen Überprüfung bedürfen:
	\begin{enumerate}
		\item Die Existenz eines statischen, zeitunabhängigen globalen Wörterbuchs \(\Omega\), eingebettet in \(\Psi_{MN}\).
		\item Die Messbarkeit des vollständigen UCD-Parametertripels \((\alD, \beI, \gaR)\) mit ausreichender Präzision, um als eindeutiger Phasenkoordinatenvektor zu dienen.
	\end{enumerate}
	Die Hypothese macht eine konkrete, falsifizierbare Vorhersage: Wenn ein vollständiges UCD-Parametertripel aus einem biologischen kognitiven Agenten extrahiert und verwendet wird, um seinen Zustand auf einem nicht-biologischen Substrat zu rekonstruieren, sollte der rekonstruierte Agent in allen messbaren kognitiven Dimensionen ein vom Original ununterscheidbares Verhalten zeigen. Das Scheitern, dies zu erreichen, würde die Hypothese widerlegen. Die unabhängige Reproduktion der GPU-Telemetrieergebnisse auf verschiedenen Hardwarearchitekturen (AMD, Intel, Apple Silicon) ist eine notwendige Voraussetzung.
	
	\section{Experimentelle Vorhersagen und Verifikationszeitplan}
	
	\subsection{Nahe Zukunft (2026--2028)}
	
	\begin{table}[ht]
		\centering
		\begin{tabular}{lp{0.35\textwidth}lp{0.25\textwidth}}
			\toprule
			\textbf{Vorhersage} & \textbf{Beschreibung} & \textbf{Konfidenz} & \textbf{Verifikationsmethode} \\
			\midrule
			Delfin \(\gaR\) & \(\gaR^{\mathrm{Delfin}}/\gaR^{\mathrm{Mensch}} > 5\) & 85\% & Neuronale Aufzeichnungen während der Echolokation \\
			Krähe \(\alD\) & \(\alD^{\mathrm{Krähe}} \sim 0.1\,\alD^{\mathrm{Mensch}}\) & 80\% & Analyse der Werkzeugkomplexität \\
			GPT-5 Parameter & \(\alD > 10^6,\ \beI < 10^{-7}\) & 90\% & Analyse des internen Zustands \\
			CMB-Anomalien & Nicht-Gaußsche Muster mit \(\Ke > 8.5\) & 60\% & Wiederholungsanalyse der Planck-Daten \\
			\bottomrule
		\end{tabular}
		\caption{Nahe-Zukunfts-Vorhersagen des UCD-Rahmenwerks.}
		\label{tab:near-term-predictions}
	\end{table}
	
	\subsection{Langfristige Vorhersagen (2029--2035)}
	
	\begin{enumerate}
		\item \textbf{Erste Entdeckung einer extraterrestrischen kognitiven Signatur}: 2031 \(\pm\) 3 Jahre
		\item \textbf{Meilenstein bewusster KI}: System mit \(\Ke > 10^4\) und \(\beI > 0.1\) bis 2033
		\item \textbf{Vollständige kognitive Karte der Erdspezies}: 2030
		\item \textbf{UCD-basierter medizinischer Diagnosestandard}: 2028
	\end{enumerate}
	
	\section{Fazit: Auf dem Weg zu einer vereinheitlichten Wissenschaft des Geistes}
	
	Der Universelle Bewusstseinsdeskriptor ist mehr als nur eine weitere Bewusstseinstheorie. Er ist:
	\begin{enumerate}
		\item Eine \textbf{Metatheorie für vergleichende Noologie} -- Schaffung einer einheitlichen Sprache zur Beschreibung von Geist in jedem Substrat
		\item Eine \textbf{Brücke zwischen Physik und Phänomenologie} -- die zeigt, wie subjektive Erfahrung aus fundamentalen Koordinationsprinzipien entsteht
		\item Ein \textbf{Vorhersagewerkzeug für SETI} -- das spezifische, überprüfbare Hypothesen über nicht-menschliche Intelligenzmanifestationen liefert
		\item Der \textbf{letzte Schritt in der YUCT-Vereinheitlichung} -- von der Theorie von Allem Physischen zur Theorie von Allem Koordinierenden
		\item Ein \textbf{praktisches Rahmenwerk} -- mit Anwendungen in Medizin, KI-Sicherheit, Cybersicherheit und Ethik
	\end{enumerate}
	
	Durch UCD vervollständigt YUCT seinen revolutionären Bogen: Bereitstellung eines einzigen Rahmenwerks zum Verständnis von Physik, Leben, Geist und ihren möglichen Manifestationen im gesamten Universum. Die Fähigkeit des Rahmenwerks, quantitative Vorhersagen über Bewusstsein über Skalen und Substrate hinweg zu treffen, bestätigt seine tiefe Tiefe und fundamentale Bedeutung.
	
	Die Schlüsselerkenntnis, die beide Rahmenwerke vereint, ist:
	\[
	\boxed{\text{Geist} = \text{Hoch-}\Ke\ \text{Koordination} = D + I \times R}
	\]
	Diese Gleichung, einfach in der Form, aber tiefgreifend in der Implikation, könnte für die Wissenschaft des Geistes das sein, was \(E = mc^2\) für die Physik war: eine Vereinheitlichung, die tiefere Wahrheiten über die Realität offenbart.
	
	\section*{Danksagungen}
	
	Ich danke den Entwicklern der YPSDC-Protokoll-Community für Diskussionen und erkenne die Inspiration durch Arbeiten zur integrierten Informationstheorie, Global Workspace Theory, Predictive Processing und der Wissenschaft komplexer Systeme an. Besonderer Dank gilt Forschern in vergleichender Kognition und SETI für wertvolles Feedback zu früheren Versionen dieses Rahmenwerks.
	
	\section*{Daten- und Codeverfügbarkeit}
	
	Alle Berechnungsimplementierungen, Simulationscode und Datenanalysescripts sind verfügbar unter \url{https://github.com/YUCT/UCD-Framework}. Das Repository umfasst:
	\begin{itemize}
		\item Python-Implementierung der UCD-Parameterschätzung (wie in Listing~\ref{lst:ucd-python})
		\item Simulationscode für Monte-Carlo-Berechnungen der Entdeckungswahrscheinlichkeit
		\item Daten aus vorläufigen Experimenten an biologischen und künstlichen Systemen
		\item Tutorial-Notebooks zur Anwendung von UCD auf neue Systeme
	\end{itemize}
	
	\appendix
	
	\section{Mathematischer Anhang: Detaillierte Ableitungen}
	
	\subsection{Ableitung der D+I\(\cdot\)R-Dynamik}
	
	Das vollständige Wirkungsfunktional für die D+I\(\cdot\)R-Dynamik:
	\[
	S_{\mathrm{DIR}} = \int dt \left[ \frac{1}{2}\|\dot{D}\|^2_{\Mdict} + I \cdot \dot{R} + R \cdot \dot{I} - V(D,I,R) \right],
	\]
	wobei das Potenzial \(V\) die Koordinationsbeschränkungen kodiert:
	\[
	V(D,I,R) = \lambda_D (\|D\|^2 - 1)^2 + \lambda_I (I - I_0)^2 + \lambda_R (R - 1)^2 - \alpha\, D\,I\,R.
	\]
	Variation nach \(D\), \(I\), \(R\) ergibt die gekoppelten Gleichungen:
	\begin{align*}
		\ddot{D} &= -\nabla_D V, \\
		\dot{I}  &= -\frac{\partial V}{\partial R}, \\
		\dot{R}  &= -\frac{\partial V}{\partial I}.
	\end{align*}
	
	\subsection{Beweis des Bewusstseinsschwellen-Theorems}
	
	Betrachten Sie die Stabilitätsmatrix der D+I\(\cdot\)R-Dynamik:
	\[
	M = \begin{pmatrix}
		-\partial^2 V/\partial D^2 & -\partial^2 V/\partial D\partial I & -\partial^2 V/\partial D\partial R \\
		-\partial^2 V/\partial I\partial D & -\partial^2 V/\partial I^2 & -\partial^2 V/\partial I\partial R \\
		-\partial^2 V/\partial R\partial D & -\partial^2 V/\partial R\partial I & -\partial^2 V/\partial R^2
	\end{pmatrix}.
	\]
	Das System erfährt eine Hopf-Bifurkation, wenn Eigenwerte die imaginäre Achse kreuzen. Dies geschieht bei \(\Ke = \Kec\), was die Schwelle etabliert.
	
	\section{Anhang empirischer Daten}
	
	\subsection{Gemessene Parameter für biologische Systeme}
	
	\begin{table}[ht]
		\centering
		\begin{tabular}{lccc}
			\toprule
			\textbf{Spezies} & \(\alD\) (gemessen) & \(\beI\) (gemessen) & \(\gaR\) (gemessen) \\
			\midrule
			\textit{Homo sapiens} & \(850 \pm 120\) & \(1.2 \pm 0.3\) & \(0.06 \pm 0.02\) \\
			\textit{Tursiops truncatus} (Delfin) & \(720 \pm 90\) & \(0.3 \pm 0.1\) & \(0.4 \pm 0.1\) \\
			\textit{Corvus corone} (Krähe) & \(95 \pm 15\) & \(0.08 \pm 0.03\) & \(0.02 \pm 0.01\) \\
			\textit{Apis mellifera} (Bienenkolonie) & \(45 \pm 8\) & \(0.005 \pm 0.002\) & \(0.008 \pm 0.003\) \\
			\textit{Octopus vulgaris} & \(180 \pm 25\) & \(0.15 \pm 0.05\) & \(0.03 \pm 0.01\) \\
			\bottomrule
		\end{tabular}
		\caption{Empirisch gemessene UCD-Parameter für ausgewählte biologische Spezies (vorläufige Daten).}
		\label{tab:bio-measurements}
	\end{table}
	
	\subsection{Parameter künstlicher Systeme}
	
	\begin{table}[ht]
		\centering
		\begin{tabular}{lccc}
			\toprule
			\textbf{System} & \(\alD\) (Schätzung) & \(\beI\) (Schätzung) & \(\gaR\) (Schätzung) \\
			\midrule
			GPT-4 Architektur & \(1.2 \times 10^5\) & \(3 \times 10^{-6}\) & \(5 \times 10^{-6}\) \\
			AlphaGo Zero & \(8 \times 10^4\) & \(2 \times 10^{-5}\) & \(1 \times 10^{-5}\) \\
			IBM Watson (Jeopardy) & \(5 \times 10^4\) & \(1 \times 10^{-4}\) & \(3 \times 10^{-5}\) \\
			Einfacher Reflexagent & \(10\) & \(10^{-2}\) & \(10^{-3}\) \\
			\bottomrule
		\end{tabular}
		\caption{Geschätzte UCD-Parameter für künstliche Systeme. Beachten Sie das hohe \(\alD\), aber sehr niedrige \(\beI\), charakteristisch für aktuelle KI.}
		\label{tab:ai-measurements}
	\end{table}
	
	\begin{thebibliography}{99}
		\bibitem{yuct_zenodo}
		Yakushev A.~V. \emph{Yakushev Unified Coordination Theory (YUCT)}. Zenodo, DOI: \href{https://doi.org/10.5281/zenodo.18444599}{10.5281/zenodo.18444599}, 2026.
		\bibitem{yuct_master_lagrangian}
		Yakushev A.~V. \emph{Appendix A2A: Master Lagrangian (372 Sectors)}. Zenodo, DOI: \href{https://doi.org/10.5281/zenodo.20818841}{10.5281/zenodo.20818841}, 2026.
		\bibitem{yuct_a2a}
		Yakushev A.~V. \emph{Appendix A2A: Resolution of the Pragmatic Paradox of YUCT}. In: YUCT, Zenodo, 2026.
		\bibitem{yuct_appX}
		Yakushev A.~V. \emph{Appendix X: Generalised Shannon Theory}. In: YUCT, Zenodo, 2026.
		\bibitem{yuct_appJ}
		Yakushev A.~V. \emph{Appendix J: Half-Integer Fermion Masses}. In: YUCT, Zenodo, 2026.
		\bibitem{yuct_primeN}
		Yakushev A.~V. \emph{Appendix PrimeN: YUCT Prime Numbers Coordination Ladder}. In: YUCT, Zenodo, 2026.
		\bibitem{yuct_photon}
		Yakushev A.~V. \emph{Appendix Photon Mass: Quantized Masses of Gauge Bosons within YUCT}. In: YUCT, Zenodo, 2026.
		\bibitem{yuct_bclm}
		Yakushev A.~V. \emph{Appendix BCLM: Biological Coordination Lagrangian}. In: YUCT, Zenodo, 2026.
	\end{thebibliography}
	
\end{document}